% Japanese introduction replaced on 2026-11-06 (README (174)); verbatim, not compiled.
\section{序論}
\label{sec:introduction}

正データ文法推論は，言語を正例のみから極限同定できるのはいつかを問う
\cite{Gold1967,Angluin1980}．
分布学習は，共通の文脈に現れる語，または語のタプルを，構成素の候補として扱う．
それは代入可能な
CFG \cite{ClarkEyraud2007,Yoshinaka2008} に始まり，多重文脈自由文法（MCFG）に拡張された．
MCFG のタプル値の非終端記号は，不連続な依存と交差依存を表す
\cite{SekiEtAl1991,Yoshinaka2009,Yoshinaka2011,Kallmeyer2010}．
Yoshinaka の有限シグネチャの枠組み \cite{Yoshinaka2015General} は，
多項式個の語からなる特徴標本による正データ同定を与える．
多項式個の語でもなお指数的に長くあり得るので，
本稿は \cite[Section~2.5.2]{EyraudHeinzYoshinaka2016} の厚データの観点に従い，有限標本 $K$ をその符号化総長
$\|K\|_+=\sum_{w\in K}(|w|+1)$ で測る．
質的な同定は
Yoshinaka の枠組みから受け継ぐ；本稿の主題は，標的の MCFG 表示に相対的な特徴標本の長さであり，
問いは次である：\emph{どの比較バイアスについてこの長さは多項式であり，
その答えはファンアウトに依存するか？}

\paragraph{設定．}
データを見る前に，ファンアウトの上限 $f$ と全射な
有限半群準同型 $s:\Sigma^+\twoheadrightarrow\mathsf S$，すなわち
\emph{型付け}を固定する．再構成は，観測された空でないタプル成分を，
それらの $s$-値が一致するときにのみ比較し，標的クラス $\Cmcf{f}{s}$
は，この制限の下でタプル代入可能な $f$-MCFL からなる（定義~\ref{def:class-cmcf}）．
すでにファンアウト2において，群型付けはこのクラスを Yoshinaka の型なしのクラスより真に大きくし，
その証拠として非文脈自由なものがあり，それらの上で型なしの再構成は過汎化する
（命題~\ref{prop:typed-separation}）．
型付けは \cite{CosteTyping2004} の意味でのクラス単位のバイアスであり，
標的ごとに選ぶものではない；他のバイアスには，制御集合によって導出を制限するもの
\cite{Takada1988,Takada1995} や，因子の比較に用いる証拠を制限するもの
\cite{CosteGaretNicolas2012,CosteNicolas2019} がある．
典型的な型付けは，
接頭辞--接尾辞窓，有限群の値，それらの積，および DFA の
遷移準同型である．群の値は
ブロックの順に合成されるので，$S_3$ のような非可換群では，
交互に並ぶ成分の順序が意味をもつ．局所群，すなわち有限半群 $\mathsf T$ であって，
$e\mathsf Te$ がすべてのべき等元 $e$ について群となるものは，
滑動因子の群寄与を加える：$\{\mathtt a,\mathtt b\}$ 上で，
$\mathtt{ab}$ の出現回数の偶奇を記録することは局所群型付けであるが，
接頭辞--接尾辞窓と群準同型の組み合わせのどれもそれを決定しない：
各窓の大きさと群準同型について，ある $m,n$ が存在して，
$\mathtt b^m\mathtt a^{n}\mathtt b^m$ と $\mathtt b^{2m}$ はその両方で一致するが，
その偶奇では異なる
（命題~\ref{prop:li-g-strict}）．窓，群およびそれらの積も局所群であるが，
DFA の遷移準同型はそうである必要はない：ある文字が
恒等写像として作用し，別の文字が少なくとも二つの状態の上の定値写像として作用するならば，その二つの変換は
$U_1$ の複製をなし，これは後の定理~\ref{thm:C} の障害である．既約な良い表示 $G$ について，
その大きさ $\|G\|$ と規則の厚さ $\bar\theta_G$ を用いる．後者は，
与えられた規則を用いる最短の文の長さのうち最大のものである．下界は，
クラス全体で特徴標本をもつすべての集合駆動型作用素について述べる；上界は，
型付きの連結をもつ Yoshinaka の再構成作用素について述べる（第~\ref{sec:learner}~節）．

\paragraph{主結果．}
アルファベット，ファンアウト $f$，規則ランクの上限 $r_{\max}\ge2$ および
型付け $s$ を固定する．

\begin{mainthm}[局所群]
\label{thm:A}
$s$ の正の像が局所群ならば，すべての $L\in\Cmcf{f}{s}$ は，
総長 $O(\|G\|^2\bar\theta_G)$ の特徴データを，すべての既約な良い表示 $G$ についてもつ
\textup(定理~\ref{thm:local-group-polydata}\textup)．
\end{mainthm}

\begin{mainthm}[線形因子]
\label{thm:B}
\begin{enumerate}[label=\textup{(\roman*)},leftmargin=*]
\item $\Sigma$ が少なくとも $4f+1$ 文字をもつか，または $f=1$ かつ $\Sigma$ が少なくとも
2文字をもつならば，すべての型付け $s$ について，
$\Cmcf{f}{s}$ の有限の標的であって，最小の既約な良い表示の定数倍以内の表示をもつものが存在し，
その特徴標本はすべて長さ
$\Omega(\|G\|\bar\theta_G)$ をもつ
\textup(命題~\ref{prop:linear-factor-all-fanout}\textup)．
\item 自明な型付けについて，特徴データに対する Yoshinaka の上界 $O(\|G\|\bar\theta_G)$
\cite{Yoshinaka2011} は定数倍を除いて達成される：
\textup{(i)} のアルファベットの条件の下で，\textup{(i)} の標的によって
\textup(系~\ref{cor:trivial-exact}\textup)．
\end{enumerate}
\end{mainthm}

\begin{mainthm}[ファンアウト1での境界]
\label{thm:C}
$f=1$ かつ $r_{\max}\ge2$ とする．
\begin{enumerate}[label=\textup{(\roman*)},leftmargin=*]
\item 有限半群の擬多様体が一様多項式特徴データ長を許すのは，
それが $\mathbf{LG}$ に含まれるとき，かつそのときに限る
\textup(定理~\ref{thm:fanout-one-data}\textup)．
\item $\mathbf{LG}$ の外では失敗は指数的である：
型付け $s_U$（二元半束 $U_1$ の上への）について，$\Cmcf{1}{s_U}$ の標的であって，
大きさと規則の厚さが $O(n)$ の既約な良い表示をもち，その
特徴標本がすべて長さ $3\cdot2^n-1$ 以上をもつものが存在する．これは，
クラス全体で特徴標本をもつすべての集合駆動型作用素について成り立つ
\textup(定理~\ref{thm:fanout-one-data}\textup)．
\end{enumerate}
\end{mainthm}

多項式時間・厚データ同定の用語
\cite{Yoshinaka2008,EyraudHeinzYoshinaka2016} では，定理~\ref{thm:A}
と~\ref{thm:C} は，局所群型付けについてクラス
$\Cmcf{f}{s}\cap\mathsf{MCFL}(f,r_{\max})$ がすべてのファンアウトでそのように同定可能であること，
およびファンアウト1では $\mathbf{LG}$ がちょうど
この性質の擬多様体の境界であることを述べている（系~\ref{cor:iptd}）．
ファンアウト1の MCFG は文脈自由文法であるので，定理~\ref{thm:C} は
文脈自由言語についての主張である；それは
姉妹原稿 \cite{kuriyamaCFG} の十分条件を特徴づけに変える．

\paragraph{残る問い．}
定理~\ref{thm:A} は，定理~\ref{thm:C}(i) の \emph{if} の方向を
すべてのファンアウトで与える．$\mathbf{LG}$ がファンアウト $f\ge2$ でもなお境界であるかどうか
は未解決である（未解決問題~\ref{op:fanout-two}）．
定理~\ref{thm:C}(ii) の証拠族はそれを決定しない：二つの穴をもつ冪の移送により，
すべての固定した型付けの下で，その成員のうち有限個を除くすべて，およびそれらの
文字が互いに素な準同型による像は，ファンアウト2でクラスを離れ，そのクラスのどの成員も
それらの長い語をその型の唯一の語として保たない
（定理~\ref{thm:doubling-no-lift}）．ファンアウト1では超多項式のデータを必要とするが
ファンアウト2では多項式のデータで足りる型付けがあれば，それは
ファンアウトを上げることが学習を真に安くし得ることを示すであろう．

\paragraph{方針．}
文法を型で細分し，その細分に型なしの評価を適用することは
高くつき得る：固定した二元型付けについて，
表示 $G_n$（$\{\mathtt a,\mathtt c\}^+$ の）であって $\|G_n\|=O(n)$ かつ
規則の厚さが高々 $2$ であり，その完全な産出型細分の厚さが少なくとも
$2^n$ であるものが存在する（命題~\ref{prop:typed-thickness-exponential}）．
したがって本稿は特徴データを元の表示を用いて評価し，
生成された証拠ではなく分布同値な代表を用いる
（系~\ref{cor:distribution-representatives}）．
所定の型をもつ短いタプルを得るために，長い部分導出を
短いものに置き換えるが，その短いものは置き換えた
部分木の群の値を保つ必要はない；すべての場合は一つの短縮原理の例である
（補題~\ref{lem:heavy-node-count}(ii)），そして語の水準での置換の段階だけが
変わる．ファンアウト1では，置き換える各部分木は一つの連続したブロックを占め，
群の値の接頭辞積に対する鳩の巣論法によって置換可能な
区間が選ばれる．ファンアウト $f\ge2$ の MCFG の導出では，部分木は最大 $f$
個のブロックを異なる根成分に与え，それらは他の部分木のブロックと交互に並ぶ．
非可換群では同時置換の効果は
この順序に依存する．ファンアウト1に対応物のない鍵となる段階は，
すべての根成分の順序付き群積を同時に保つ置換の
\emph{所有者ごとの同時選択}である
（命題~\ref{prop:block-replacement-interface}）．
局所群型付けでは，群の値は滑動因子にわたる積であり，滑動因子は
ブロックの境界をまたぐ；ブロックが互いに遠く離れている所有者を選び，
各ブロックを変化しない境界の文字で広げることにより，この場合は
同じ選択に帰着する（補題~\ref{lem:lg-buffered-owners}）．下界は，
その部分集合がすべてクラスにとどまる標識族
（命題~\ref{prop:linear-factor-all-fanout}）と稀な型
（命題~\ref{prop:rare-type}）を用いる．ファンアウト1に対応物のない
第二の要素である二つの穴をもつ冪の移送は，ファンアウト1の指数的な
障害をファンアウト2で排除する（補題~\ref{lem:power-transfer}，
定理~\ref{thm:doubling-no-lift}）．

\paragraph{先行研究との関係．}
Yoshinaka は $(k,\ell)$-代入可能な CFG の特徴データを
厚さのパラメータで評価し \cite{Yoshinaka2008}，多次元
代入可能な MCFG を正データから学習し \cite{Yoshinaka2009,Yoshinaka2011}，
有限の証拠の個数をもつ有限シグネチャの学習原理を与えた
\cite{Yoshinaka2015General}；質的な結果と，ここで用いる
Const/Link/Comp の構成は彼のものである．有限半群の
構造理論と局所群の表現は古典的である
\cite{Straubing1985,Tilson1987,AlmeidaEtAl2009}．本稿の寄与は，
符号化長の解析と上記の境界の結果である．他の
弱文脈依存の形式体系の正データからの学習は
\cite{OatesEtAl2006} にあり，質問や双対文脈集合を用いた MCFG の学習は
\cite{Yoshinaka2010,ClarkYoshinaka2014,KanazawaYoshinaka2023} にある；
\cite{DuncanEtAl2026Substitution} の MCFL の
代入補題は所属の条件であり，学習の基準ではない．

\paragraph{姉妹 CFG 原稿との関係．}
姉妹原稿 \cite{kuriyamaCFG} はファンアウト1を扱う．要するに，それは
$\mathbf{LI}\vee\mathbf{Grp}$ の型付けについての多項式データを
十分条件として，また $U_1$ の上への一つの型付けについての指数的データを証明する；ここでは
ファンアウト1の境界をちょうど（$\mathbf{LG}$）決定し，
$\mathbf{LG}$ に対する上界と線形因子の下界はすべてのファンアウトで成り立つ．
表~\ref{tab:companion-overlap} はその結果を本稿の結果から区別する（第2列の番号は
姉妹原稿のものを指す）．姉妹原稿のどの結果も前提としては用いない：
その括弧付き倍増族
\cite[Proposition~6.4]{kuriyamaCFG} の代わりに，三文字の変種を用い，それは
付録~\ref{app:bracketed-family} で証明する．また，その表現力についての結果（カウンタと
Dyck の障害，合同文法との比較）は用いない．

\begin{table}[htbp]
\caption{姉妹 CFG 原稿 \cite{kuriyamaCFG}
（ファンアウト1）の結果と本稿の結果．}
\label{tab:companion-overlap}
\centering
\footnotesize
\setlength{\tabcolsep}{3pt}
\renewcommand{\arraystretch}{1.15}
\begin{tabularx}{\linewidth}{@{}>{\raggedright\arraybackslash}p{0.27\linewidth}>{\raggedright\arraybackslash}p{0.25\linewidth}>{\raggedright\arraybackslash}X@{}}
\hline
主題 & 姉妹原稿 \cite{kuriyamaCFG} & 本稿\\
\hline
型付き再構成，Gold 学習器，型付きの厚さによる証拠の上界
 & 定理~4.2，定理~5.2--5.3，系~5.4；補題~6.2，系~6.3
 & タプル版（第~\ref{sec:learner}~節，
   定理~\ref{thm:data-from-typed-thickness}）；質的な結果は
   Yoshinaka に帰する；新しくない\\
窓および窓--群の上界
 & 命題~3.5，補題~7.1，定理~7.2，補題~7.4，系~7.5
 & $f=1$ では新しくない；$f\ge2$ では新しい：交互に並ぶ
   成分にわたる所有者ごとの置換（命題~\ref{prop:block-replacement-interface}）\\
局所群型付け
 & 扱わない（$\mathbf{LI}\vee\mathbf{Grp}$ のみ；補題~7.6）
 & データ長について新しい，すべての $f$ で（定理~\ref{thm:lg-characterization}
   と~\ref{thm:local-group-polydata}）；$\mathbf{LG}$ の $U_1$ による特徴づけと
   半群の代数は古典的である \cite{AlmeidaEtAl2009}\\
$U_1$ の上への型付けの下での倍増族
 & 命題~6.4，系~6.6，付録~C（五文字，括弧）
 & 同じ上界・下界をもつ三文字の変種，ここで証明する
   （命題~\ref{prop:three-letter-family}）\\
稀な型による下界
 & 命題~6.4(iii)--(iv)，一つの語について
 & すべての実現型とすべての $f$（命題~\ref{prop:rare-type}）\\
ファンアウト1での擬多様体の境界
 & $\mathbf{LI}\vee\mathbf{Grp}$ に対する十分性のみ（注意~7.7）
 & 新しい（定理~\ref{thm:fanout-one-data}）\\
型付きのクラスと型なしのクラス
 & $f=1$ での分離
 & $f\ge2$ で，Yoshinaka のクラスの外にある非文脈自由な成員
   （命題~\ref{prop:typed-separation}）\\
因子 $\|G\|\bar\theta_G$ の必要性
 & ---
 & すべての $f$ で新しい（命題~\ref{prop:linear-factor-all-fanout}）\\
ファンアウト2：冪の移送
 & ---
 & 新しい（定理~\ref{thm:doubling-no-lift}）\\
\hline
\end{tabularx}
\end{table}

\paragraph{主張の範囲．}
すべての上界・下界は，指定した集合駆動型作用素に対する特徴標本の符号化総長に関するものであり，
述べた表示に相対的である；
それらは文法の固有の大きさも Gold の
収束の段階も評価せず，その指数については
第~\ref{sec:conclusion}~節の終わりで論じる．データの下界は
判別標的のみから得られる（命題~\ref{prop:rare-type}
と~\ref{prop:linear-factor-all-fanout}）；厚い細分
（命題~\ref{prop:typed-thickness-exponential}）は証拠の
構成を制限するのであって，データを制限するのではない（系~\ref{cor:distribution-representatives}）．
型付けは標的を見る前に固定する；それを細分するとクラスは大きくなるが，
再構成作用素が変わるので，本稿はそれを推論も最適化もしない．
表~\ref{tab:scope} は，各主張について，
固定するパラメータ，それが関わる作用素，およびその付帯条件を挙げる；上界の定数は
固定したパラメータのみに依存する．

\begin{table}[htbp]
\caption{主な主張の範囲．「標準」は
$\mathcal B^{\le r_{\max}}_{f,s}$ と $\mathcal B_{f,s}$ の両方（一つの共通の標本で）を意味する；
「クラス全体」は，クラス全体で特徴標本をもつ
すべての集合駆動型作用素を意味する．$G$ は $L\cap\Sigma^+$ の既約な良い表示であって，
ファンアウト $\le f$ かつランク $\le r_{\max}$ のものである．}
\label{tab:scope}
\centering
\footnotesize
\setlength{\tabcolsep}{3pt}
\renewcommand{\arraystretch}{1.15}
\begin{tabularx}{\linewidth}{@{}>{\raggedright\arraybackslash}p{0.2\linewidth}>{\raggedright\arraybackslash}p{0.22\linewidth}>{\raggedright\arraybackslash}p{0.17\linewidth}>{\raggedright\arraybackslash}X@{}}
\hline
主張 & 固定するもの & 作用素 & 評価と条件\\
\hline
定理~\ref{thm:A}（定理~\ref{thm:local-group-polydata}）
 & $f$，$r_{\max}$，$\Sigma$，$s$ で $s(\Sigma^+)\in\mathbf{LG}$ のもの
 & 標準
 & 上界 $O(\|G\|^2\bar\theta_G)$，すべての $L\in\Cmcf{f}{s}$ とすべての $G$ について\\
定理~\ref{thm:B}(i)（命題~\ref{prop:linear-factor-all-fanout}）
 & $f$，任意の $s$
 & 標準（$r\ge1$），クラス全体
 & 有限の標的に対する下界 $\Omega(\|G\|\bar\theta_G)$；$|\Sigma|\ge4f+1$（$|\Sigma|\ge2$，ただし $f=1$ の場合）を要する；
   ランク2の表示\\
定理~\ref{thm:B}(ii)（系~\ref{cor:trivial-exact}）
 & $f$，$r_{\max}\ge1$，$s_{\mathrm{triv}}$
 & $\mathcal B_{f,s_{\mathrm{triv}}}$
 & 上界 $O(\|G\|\bar\theta_G)$（Yoshinaka \cite{Yoshinaka2011}）；$r_{\max}\ge2$ かつ $|\Sigma|\ge4f+1$（$\ge2$，ただし $f=1$ の場合）のとき達成される\\
定理~\ref{thm:C}(i)（定理~\ref{thm:fanout-one-data}）
 & $f=1$，$r_{\max}\ge2$，一つの擬多様体
 & $\mathcal B^{\le r_{\max}}_{1,s}$（定義~\ref{def:uniform-polydata}）
 & 一様多項式データをもつのは $\mathbf V\subseteq\mathbf{LG}$ のとき，かつそのときに限る\\
定理~\ref{thm:C}(ii)
  & $f=1$，三文字上の $s_U$
 & 標準，クラス全体
 & 下界 $3\cdot2^n-1$，大きさと厚さが $O(n)$ の表示に対して\\
系~\ref{cor:iptd}
 & $f$，$r_{\max}$，$s$
 & 集合駆動型
 & $\mathbf{LG}$ について IPTtD，すべての $f$ で；$f=1$ ではちょうど $\mathbf{V}\subseteq\mathbf{LG}$ について\\
定理~\ref{thm:doubling-no-lift}
 & 任意の $f\ge2$，任意の $s$
 & なし（所属）
  & 定理~\ref{thm:C}(ii) の族とその文字が互いに素な像は $\Cmcf{f}{s}$ を離れ，その強制された語はそこで決してその型の唯一の語ではない；$f\ge2$ でのデータの評価はない\\
未解決
 & 
 & 
 & $\mathbf{LG}$ が $f\ge2$ で境界であるかどうか（未解決問題~\ref{op:fanout-two}）\\
\hline
\end{tabularx}
\end{table}

\paragraph{構成．}
第~\ref{sec:prelim}~節は設定を固定し，第~\ref{sec:learner}~節は
再構成作用素とその質的な保証を思い出し，
型付きのクラスを型なしのクラスから分離する（第~\ref{subsec:typed-separation}~節）．
第~\ref{sec:char-data}~節は特徴データを型付きの厚さと
分布の代表に帰着させる．第~\ref{sec:upper}~節は
定理~\ref{thm:A} を証明し，第~\ref{sec:linear}~節は定理~\ref{thm:B} を証明し，
第~\ref{sec:fanout-one}~節は定理~\ref{thm:C} を証明し，
その境界を多項式時間・厚データ同定の言葉で言い直す．
第~\ref{sec:fanout-two}~節は証拠族がファンアウト2に持ち越されないことを示し，
第~\ref{sec:conclusion}~節は指数と
未解決問題を論じる．付録は技術的な証明を含む．

